% Source: https://artofproblemsolving.com/community/c3h1656183_euler_would_be_proud_paul_msa116
% Stable identity: aops-topic-1656183-q01
A regular octahedron consists of two square pyramids glued along their square bases.

\begin{asy}
usepackage("amsmath");
size(3.5cm);
defaultpen(fontsize(10pt));
pair T=(0,2.2),Bo=(0,-2.2);pair a=(-1.8,-0.25),c=(1.8,0.25),bb=(-0.6,0.55),d=(0.6,-0.55);
draw(T--a--Bo--c--cycle,linewidth(0.9));draw(T--d--Bo,linewidth(0.9));draw(a--d--c,linewidth(0.9));
draw(T--bb,dashed);draw(Bo--bb,dashed);draw(a--bb--c,dashed);
dot(T);dot(Bo);dot(a);dot(c);dot(d);dot(bb);
label("\mbox{}",(min(currentpicture,user=true)+max(currentpicture,user=true))/2,black); // keeps the text after the figure black when LaTeX typesets the labels inline
\end{asy}

It has $F=8$ triangular faces and $V=6$ vertices (the $4$ corners of the square and the $2$ apexes). Each face has $3$ edges and each edge lies on $2$ faces, so $E=\dfrac{8\cdot3}{2}=12$; Euler's formula $V-E+F=6-12+8=2$ confirms this. The sum is $8+6+12=\boxed{\textbf{(f)}\ 26}$.
